\scanpage{260}

le sous-groupe $Z$ de $T_G$ comme une sorte de sous-groupe diagonal, avec un diagramme commutatif

\diagram{\begin{tikzcd}[ampersand replacement=\&,column sep=small]
\& 1\arrow[r] \& \unsure{\mathfrak{z}|\mathfrak{Z}}\arrow[r]\arrow[d,equal] \& \widetilde Z\arrow[r,"c"]\arrow[d,"\mathrm{diag}"] \& Z\arrow[r,""{name=nblink,coordinate}]\arrow[d,hook,"\mathrm{incl}."] \& 1 \& |[alias=nbtip]| {} \arrow[from=nblink,to=nbtip,bend left=43]\\
(130)\&1\arrow[r] \& \mathfrak{z}\arrow[r,swap,"{z\mapsto(z,z^{-1})}"] \& \widetilde Z\times\widetilde Z\arrow[r] \& T_G\arrow[r] \&1
\end{tikzcd}}

\margin{NB on a supposé $c$ surjectif (cas où la ``condition d'isomorphie'' est vérifiée)}

\deleted{déjà le} \deleted{mais ce diagramme n'est commutatif,}

l'inclusion $Z\subset T_G$ étant donc déduite (ou mieux sur le sous-groupe $c(\widetilde Z)$ de $Z$) par passage au quotient à partir de l'inclusion diagonale $\widetilde Z\longrightarrow\widetilde Z\times\widetilde Z$. Mais pour que cette inclusion \deleted{\unsure{transforme}} \inserted{applique bien} $\mathfrak{z}=\operatorname{Ker} c$ dans $\mathfrak{z}$ \deleted{$\mathfrak{z}=$} $\operatorname{Ker}(\widetilde Z\times\widetilde Z\longrightarrow T_G)$, il faut (et suffit) qu'on ait

\[z^2=1\qquad\forall z\in\mathfrak{z},\tag*{(131)}\]

et alors le premier carré de (130) est bien commutatif (l'hom $\mathfrak{z}\longrightarrow\widetilde Z\times\widetilde Z$ de la deuxième ligne \deleted{\unsure{s'identifie à}} s'écrit aussi, alors $z\unsure{-|\mapsto}(z,z)$ i.e. est \deleted{\unsure{induit par}} \inserted{le composé}

\[\mathfrak{z}\overset{\mathrm{diag}}{\unsure{-|\longrightarrow}}\mathfrak{z}\times\mathfrak{z}
\underset{\substack{\mathrm{incl}.\\\mathrm{can}}}{\lhook\joinrel\longrightarrow}
\widetilde Z\times\widetilde Z\quad).\]

La commutativité du deuxième carré de (130) s'écrit sous la forme

\[c(\widetilde\lambda)=u_0(\widetilde\lambda)u_1(\widetilde\lambda)\qquad\forall\widetilde\lambda\in\widetilde Z\tag*{(132)}\]

(qui \deleted{\unsure{imp}} pour $\deleted{\unsure{z}}\widetilde\lambda=z\in\mathfrak{z}$ se réduit à (131)). On voit immédiatement que les deux mbes de (132) \deleted{\unsure{appliq}} ont $=$ effet sur $(s_0,t_0)$

\[(s_0,t_0)\longmapsto(\lambda s_0,\lambda t_0),\quad\text{où }\lambda=c(\widetilde\lambda),\]

% archive sheet number: 259
